\section{Vzorová řešení z posledních zkoušek}

{\small Výběr je omezený na společné okruhy a specializaci UI-SU / UI-SU, UI-ZPJ, které jsou
pokryté v těchto poznámkách. Každé řešení je psané jako modelová odpověď na tříbodovou otázku:
nejprve definice, potom konstrukce nebo výpočet, nakonec hraniční případ nebo praktický komentář.}

\subsection{2026-jaro, otázka 1: DFA pro porovnání modulárních hashů}

Zadání: pro prvočíslo $p$ je $h(w)$ hodnota binárního slova $w$ modulo $p$. Sestrojte DFA pro
\[
L_p=\{u\#v\mid u,v\in\{0,1\}^*,\ h(u)=h(v)\}.
\]

\begin{odpoved}
\textbf{Definice.} Deterministický konečný automat je pětice
$A=(Q,\Sigma,\delta,q_0,F)$, kde $Q$ je konečná neprázdná množina stavů,
$\Sigma$ konečná neprázdná abeceda, $\delta:Q\times\Sigma\to Q$ přechodová funkce,
$q_0\in Q$ počáteční stav a $F\subseteq Q$ množina přijímajících stavů. Pokud připustíme
parciální $\delta$, chybějící přechody lze doplnit mrtvým stavem.

\textbf{Konstrukce.} Označme $[p]=\{0,\ldots,p-1\}$. Automat má
\[
Q=[p]\cup([p]\times[p]),\qquad \Sigma=\{0,1,\#\},\qquad q_0=0,
\]
\[
F=\{(i,i)\mid i\in[p]\}.
\]
Stav $i$ znamená, že jsme stále před znakem $\#$ a hodnota už přečteného prefixu $u'$ je
$h(u')=i$. Stav $(i,j)$ znamená, že jsme za znakem $\#$, pamatujeme $i=h(u)$ a pro dosavadní
prefix $v'$ platí $j=h(v')$.

Přechody jsou pro $i,j\in[p]$:
\[
\delta(i,0)=2i\bmod p,\qquad
\delta(i,1)=(2i+1)\bmod p,\qquad
\delta(i,\#)=(i,0),
\]
\[
\delta((i,j),0)=(i,2j\bmod p),\qquad
\delta((i,j),1)=(i,(2j+1)\bmod p).
\]
Přechod na druhý znak $\#$ je nedefinovaný, případně vede do žumpy.

\textbf{Příklad pro $p=2$.} Před $\#$ stačí rozlišovat paritu hodnoty. Ze stavu $0$ vede $0$ do
$0$, $1$ do $1$ a $\#$ do $(0,0)$. Ze stavu $1$ vede $0$ do $0$, $1$ do $1$ a $\#$ do $(1,0)$.
Za $\#$ se první složka nemění a druhá se aktualizuje stejně modulo $2$:
z každého $(i,j)$ vede $0$ do $(i,0)$ a $1$ do $(i,1)$. Přijímající jsou $(0,0)$ a $(1,1)$.
Tím automat přijme přesně slova s jedním oddělovačem a shodnými zbytky.
\end{odpoved}

\subsection{2026-jaro, otázka 27: CSP a AC-3}

Zadání: definujte binární CSP a hranovou konzistenci, rozhodněte vztah ke splnitelnosti a
aplikujte AC-3 na domény
\[
D_A=\{1\},\quad D_B=D_C=D_D=\{1,2,3\}
\]
s podmínkami $C+D=4$, $A+B\le 3$, $B+C\ge 4$.

\begin{odpoved}
\textbf{Definice.} Binární CSP je dáno konečnou množinou proměnných, konečnou doménou
$D_X$ pro každou proměnnou $X$ a množinou binárních podmínek, tedy relací nad dvojicemi
proměnných. Výstupem je úplné přiřazení hodnot všem proměnným splňující všechny podmínky,
nebo zjištění, že žádné neexistuje. Orientovaná hrana $(X,Y)$ je hranově konzistentní, pokud
pro každé $x\in D_X$ existuje $y\in D_Y$ takové, že dvojice splní podmínku mezi $X$ a $Y$.

\textbf{Splnitelnost.} Hranová konzistence negarantuje splnitelnost. Protipříklad: tři
proměnné $X_1,X_2,X_3$ s doménou $\{0,1\}$ a podmínkami $X_i\ne X_j$ pro každou dvojici.
Každá dvojice má podporu, ale tři proměnné dvěma hodnotami obarvit nelze. Naopak hranově
nekonzistentní problém může být splnitelný, protože AC-3 může teprve odstranit hodnoty,
které v žádném řešení nejsou.

\textbf{AC-3 na zadaném příkladu.} Začneme všemi orientovanými hranami. Podmínka
$A+B\le3$ při $D_A=\{1\}$ dává pro $B$ podporu jen hodnotám $1,2$, protože $1+3>3$.
Tedy $D_B$ se zmenší na $\{1,2\}$.

Podmínka $B+C\ge4$ teď čistí obě strany. Pro $B=1$ je potřeba $C=3$, pro $B=2$ stačí
$C=2$ nebo $3$, takže $B$ zůstává $\{1,2\}$. Hodnota $C=1$ nemá podporu v $B=\{1,2\}$,
protože $1+1<4$ i $2+1<4$. Tedy $D_C=\{2,3\}$.

Podmínka $C+D=4$ s $D_C=\{2,3\}$ ponechá $D_D=\{1,2\}$: pro $C=2$ je podpora $D=2$,
pro $C=3$ je podpora $D=1$; hodnota $D=3$ by vyžadovala $C=1$, která už v doméně není.
Zpětná kontrola nezmění $C$, protože $C=2$ má $D=2$ a $C=3$ má $D=1$.

Výsledné hranově konzistentní domény jsou
\[
D_A=\{1\},\qquad D_B=\{1,2\},\qquad D_C=\{2,3\},\qquad D_D=\{1,2\}.
\]
Například přiřazení $A=1,B=1,C=3,D=1$ ukazuje, že tento konkrétní problém je splnitelný.
\end{odpoved}

\subsection{2026-jaro, otázka 28: Shlukování pomocí k-means}

Zadání: popište k-means, jeho inicializaci, optimalizovanou ztrátu, konvergenci a praktické
důsledky.

\begin{odpoved}
\textbf{Účel a algoritmus.} k-means je metoda učení bez učitele pro shlukování bodů do $K$
skupin podle geometrické podobnosti. Na vstupu jsou body $x_1,\ldots,x_N$ a počet shluků
$K$. Algoritmus:
\begin{enumerate}
  \item Inicializuj centroidy $\mu_1,\ldots,\mu_K$, například náhodným výběrem bodů.
  \item Přiřaď každý bod k nejbližšímu centroidu.
  \item Každý centroid nahraď průměrem bodů přiřazených do jeho shluku.
  \item Opakuj přiřazení a přepočet centroidů, dokud se nezmění přiřazení nebo centroidy.
\end{enumerate}
Inicializace kmeans++ zvolí první centroid náhodně a další vybírá s pravděpodobností
úměrnou druhé mocnině vzdálenosti od nejbližšího už zvoleného centroidu.

\textbf{Ztrátová funkce.} Minimalizuje se součet čtverců vzdáleností bodů od centroidů:
\[
J=\sum_{k=1}^K\sum_{x_i\in C_k}\|x_i-\mu_k\|^2.
\]
Při pevných centroidech je nejlepší přiřazení k nejbližšímu centroidu. Při pevném
přiřazení je nejlepší centroid aritmetický průměr bodů ve shluku.

\textbf{Konvergence a důsledky.} Každý krok algoritmu hodnotu $J$ nezvětší: přiřazení k
nejbližšímu centroidu ji sníží nebo ponechá, přepočet průměru je optimální centroid pro
daný shluk. Ztráta je omezená zdola nulou a počet možných přiřazení bodů do $K$ shluků je
konečný, takže algoritmus skončí v lokálním optimu. Negarantuje globální optimum.

Prakticky to znamená, že výsledek závisí na inicializaci. Běžně se pouští více běhů s
různými počátečními centroidy a vybere se řešení s nejmenší ztrátou. Metoda předpokládá
spíše kompaktní kulové shluky a je citlivá na měřítko příznaků, proto je často potřeba
standardizace.
\end{odpoved}

\subsection{2026-jaro, otázka 29: Kombinace modelů}

Zadání: vysvětlete bagging a boosting a rozhodněte, kdy většinové hlasování modelů s
accuracy $0.70$, $0.75$ a $0.80$ překoná nejlepší jednotlivý model.

\begin{odpoved}
\textbf{Bagging a boosting.} Bagging trénuje více modelů nezávisle na různých bootstrap
vzorcích trénovací množiny, tedy na výběrech s opakováním. Výstupy se agregují
hlasováním nebo průměrem. Typický příklad je random forest, kde se navíc v každém uzlu
stromu náhodně omezuje množina uvažovaných příznaků.

Boosting trénuje modely sekvenčně. Další model klade větší důraz na příklady, které
předchozí modely klasifikovaly špatně. Příklady jsou AdaBoost, gradient boosting a
XGBoost.

\textbf{Hlasování tří modelů.} Ensemble ze tří binárních klasifikátorů chybuje právě tehdy,
když na dané instanci chybují aspoň dva modely. Proto nestačí znát jen samostatné accuracy.
Rozhodující je překryv množin chyb.

Pokud jsou chyby modelů co nejvíc komplementární, může ensemble překonat nejlepší model.
V extrémním ideálním případě jsou množiny chyb tří modelů disjunktní. Potom na každém
vzorku chybuje nejvýše jeden model, většina hlasuje správně a accuracy ensemblu je
$100\,\%$.

Naopak pokud se chyby silně překrývají, ensemble může být horší než nejlepší model. Jestliže
chyby modelu $M_2$ leží uvnitř chyb modelu $M_1$, hlasování na těchto instancích nepomůže
a accuracy může být nejvýše $0.75$. Pokud navíc chyby modelu $M_3$ pokryjí zbytek chyb
modelu $M_1$, na každé chybě $M_1$ chybují aspoň dva modely a ensemble může mít accuracy
jen $0.70$.

\textbf{Hraniční pointa.} Ensemble zlepšuje výsledek hlavně tehdy, když jsou modely dostatečně
přesné a jejich chyby nejsou silně korelované. Samotná vyšší accuracy jednotlivých modelů
bez diverzity chyb nezaručuje lepší většinové hlasování.
\end{odpoved}

\subsection{2026-jaro, otázka 30: Rozhodovací strom}

Zadání: popište stavbu rozhodovacího stromu, s kořenem v atributu $D$ určete další dělení
pro zadaná data a vysvětlete prořezávání.

\begin{odpoved}
\textbf{Algoritmus.} Rozhodovací strom stavíme hladově. V uzlu máme podmnožinu trénovacích
dat. Pokud je čistá, vytvoříme list s danou třídou. Jinak vybereme test atributu, který
nejvíc zlepší čistotu potomků, typicky maximalizuje entropický zisk nebo sníží Gini
nečistotu. Entropie dvoutřídního uzlu s počty $a,b$ je
\[
H(a,b)=-\frac{a}{a+b}\log_2\frac{a}{a+b}
-\frac{b}{a+b}\log_2\frac{b}{a+b}.
\]
Entropický zisk je entropie před dělením minus vážený průměr entropií po dělení.

\textbf{Kořen $D$.} V zadaných datech má $D$ hodnoty $3$ a $6$, takže přirozený test je
$D\le4.5$. Pro $D=3$ jsou řádky $2,6,7$ s třídami $1,0,1$. Pro $D=6$ jsou řádky
$0,1,3,4,5$ s třídami $0,0,1,1,0$.

V pravém podstromu $D=6$ je podle nástinu třeba porovnat atributy $A,B,C$. Dělení podle
$A$ nebo $B$ dává váženou entropii přibližně
\[
\frac35H(1,2)+\frac25H(1,1)\approx 0.95.
\]
Dělení podle $C$ dává
\[
\frac15H(1,0)+\frac45H(3,1)\approx 0.65.
\]
Proto je v pravém podstromu nejlepší dělit podle $C$. Větev $C=6$ je čistý list $Y=1$,
větev $C=8$ obsahuje tři nuly a jednu jedničku a pokračuje dalším testem.

V levém podstromu $D=3$ lze oddělit řádek $6$ pomocí $C=6$, takže $C=6$ dá list $Y=0$ a
$C=8$ dá list $Y=1$. Neprořezaný strom do maximální hloubky tedy lze popsat pravidly:
\[
\begin{array}{ll}
D=3,\ C=6 &\Rightarrow Y=0,\\
D=3,\ C=8 &\Rightarrow Y=1,\\
D=6,\ C=6 &\Rightarrow Y=1,\\
D=6,\ C=8,\ A\le2 &\Rightarrow Y=0,\\
D=6,\ C=8,\ A>2 &\Rightarrow Y=0\ \text{s jednou trénovací chybou, případně se dále štěpí.}
\end{array}
\]
Poslední uzel je právě místo, kde oficiální nástin uvádí prořezání.

\textbf{Prořezávání.} Prořezávání omezuje přeučení: složitý strom může mít malou trénovací
chybu, ale horší chybu na validačních nebo testovacích datech. V tomto příkladu se uzel
$A\le2.0$ nahradí listem, protože se tím nezhorší klasifikační chyba a strom má méně listů.
Obecné cost-complexity pruning používá kritérium
\[
\alpha_{\mathrm{eff}}=\frac{R(t)-R(T_t)}{|T_t|-1},
\]
kde $R(T_t)$ je chyba podstromu, $R(t)$ chyba po nahrazení listem a $|T_t|$ počet listů
podstromu. Postupně se odstraňují větve s nejmenším efektivním zlepšením pod zvoleným
prahem.
\end{odpoved}

\subsection{2025-podzim, otázka 1: Třídy složitosti}

Zadání: definujte $P$, $NP$, polynomiální redukci, NP-úplnost a rozhodněte, zda je v $NP$
problém faktorizace čísla $x$ na dva stejně dlouhé netriviální desítkové faktory.

\begin{odpoved}
\textbf{Definice.} Třída $P$ obsahuje rozhodovací problémy, které lze rozhodnout
deterministickým algoritmem v čase polynomiálním ve velikosti vstupu. Třída $NP$ obsahuje
rozhodovací problémy, pro něž existuje polynomiálně dlouhý certifikát a deterministický
verifikátor běžící v polynomiálním čase:
\[
x\in L\iff \exists y,\ |y|\le p(|x|): V(x,y)=1.
\]
Polynomiální redukce $A\le_P B$ je funkce $f$ spočitatelná v polynomiálním čase taková, že
$x\in A$ právě tehdy, když $f(x)\in B$. Problém je NP-těžký, pokud se na něj polynomiálně
redukuje každý problém z $NP$. Je NP-úplný, pokud je NP-těžký a zároveň leží v $NP$.

\textbf{Daný problém je v $NP$.} Vstupem je desítkový zápis přirozeného čísla $x$. Jako
certifikát stačí číslo $a$. Verifikátor zkontroluje:
\begin{enumerate}
  \item $1<a<x$,
  \item zda $a$ dělí $x$ beze zbytku,
  \item položí $b=x/a$ a ověří $1<b<x$,
  \item porovná délky desítkových zápisů $a$ a $b$ bez počátečních nul.
\end{enumerate}
Pokud všechny testy projdou, přijme.

\textbf{Polynomiálnost.} Certifikát $a$ má nejvýše tolik číslic jako $x$, tedy délku
lineární ve velikosti vstupu. Dělení velkých čísel, porovnání a zjištění délky zápisu lze
provést v polynomiálním čase vzhledem k počtu číslic $x$. Proto problém leží v $NP$.
Netvrdíme tím, že je NP-úplný, ani že neumíme faktorizovat obecně v polynomiálním čase.
\end{odpoved}

\subsection{2025-podzim, otázka 2: Čtení binárního souboru}

Zadání: navrhněte C\# třídy `Reader` a `Node` pro binární soubor s hlavičkou typů atributů,
big-endian čísly a uzly s délkou dat. `Node` má držet surová data a umožnit konstantní
lokalizaci atributu; implementujte `getUInt16` a `getSInt16`.

\begin{odpoved}
\textbf{Datový návrh.} `Reader` drží odkaz na `BinaryFile`, počet atributů a pole typů.
Hlavička má dva bajty počtu atributů a potom jeden bajt na typ. Pozice kořenového uzlu je
tedy $2+\texttt{attributes}$. `readNode` na dané pozici přečte dvoubajtovou délku dat uzlu,
potom přesně tolik bajtů dat a vytvoří `Node`.

`Node` ukládá data uzlu jako `byte[]` přesně ve formátu souboru. Navíc si v konstruktoru
jedním průchodem spočítá offset začátku každého atributu v datech. Proto lze i-tý atribut
lokalizovat v čase $O(1)$. U řetězce se posun zvětší o $2+\text{délka}$, u čísel o $2$ a u
linku o $8$.

\begin{verbatim}
enum AttrType { TUInt16 = 0x48, TLink = 0x4C,
                TString = 0x53, TSInt16 = 0x68 }

class Reader {
    private readonly BinaryFile _file;
    private readonly AttrType[] _types;
    private readonly long _root;

    public Reader(BinaryFile file) {
        _file = file;

        byte[] countBuf = new byte[2];
        _file.readBytes(0, countBuf, 2);
        int count = (countBuf[0] << 8) | countBuf[1];

        byte[] typeBuf = new byte[count];
        _file.readBytes(2, typeBuf, count);

        _types = new AttrType[count];
        for (int i = 0; i < count; i++)
            _types[i] = (AttrType)typeBuf[i];

        _root = 2 + count;
    }

    public int attributes() { return _types.Length; }
    public AttrType getType(int i) { return _types[i]; }
    public long getRoot() { return _root; }

    public Node readNode(long filepos) {
        byte[] lenBuf = new byte[2];
        _file.readBytes(filepos, lenBuf, 2);
        int len = (lenBuf[0] << 8) | lenBuf[1];

        byte[] data = new byte[len];
        _file.readBytes(filepos + 2, data, len);
        return new Node(_types, data);
    }
}

class Node {
    private readonly AttrType[] _types;
    private readonly byte[] _data;
    private readonly int[] _offsets;

    public Node(AttrType[] types, byte[] data) {
        _types = types;
        _data = data;
        _offsets = new int[types.Length];

        int pos = 0;
        for (int i = 0; i < types.Length; i++) {
            _offsets[i] = pos;
            switch (types[i]) {
                case AttrType.TUInt16:
                case AttrType.TSInt16:
                    pos += 2;
                    break;
                case AttrType.TLink:
                    pos += 8;
                    break;
                case AttrType.TString:
                    int len = (_data[pos] << 8) | _data[pos + 1];
                    pos += 2 + len;
                    break;
            }
        }
    }

    public int getUInt16(int i) {
        int p = _offsets[i];
        return (_data[p] << 8) | _data[p + 1];
    }

    public int getSInt16(int i) {
        int u = getUInt16(i);
        if ((u & 0x8000) != 0) return u - 0x10000;
        return u;
    }
}
\end{verbatim}

\textbf{Hraniční případy.} Big-endian znamená, že první bajt je vyšší část čísla. Pro
znaménkové šestnáctibitové číslo ve dvojkovém doplňku mají hodnoty s nastaveným bitem
$0x8000$ záporný význam, proto od unsigned hodnoty odečteme $0x10000$. Metody `get...`
nemají kontrolovat typ atributu, protože zadání tuto odpovědnost dává volajícímu.
\end{odpoved}

\subsection{2025-podzim, otázka 28: Evaluace binárního klasifikátoru}

Zadání: na 10000 instancích je 1 procento anomálií. První nástroj má accuracy $99.5\,\%$
a recall $90\,\%$. Druhý má recall $81\,\%$ a jen $10\,\%$ nahlášených anomálií je falešně
pozitivních. Spočítejte matice konfuze a metriky.

\begin{odpoved}
V datech je $100$ pozitivních instancí a $9900$ negativních.

\textbf{První nástroj.} Recall $0.90$ znamená
\[
TP=90,\qquad FN=10.
\]
Accuracy $0.995$ na $10000$ instancích znamená $9950$ správných predikcí. Proto
\[
TN=9950-90=9860,\qquad FP=9900-9860=40.
\]
Metriky:
\[
\mathrm{precision}=\frac{TP}{TP+FP}=\frac{90}{130}\approx0.69,
\]
\[
\mathrm{recall}=0.90,\qquad
\mathrm{accuracy}=0.995,
\]
\[
F_1=\frac{2PR}{P+R}\approx\frac{2\cdot0.69\cdot0.90}{0.69+0.90}\approx0.78.
\]

\textbf{Druhý nástroj.} Recall $0.81$ znamená
\[
TP=81,\qquad FN=19.
\]
Jen $10\,\%$ hlášení je falešně pozitivních, tedy precision je $0.90$:
\[
0.90=\frac{TP}{TP+FP}=\frac{81}{81+FP}.
\]
Z toho $FP=9$ a
\[
TN=9900-9=9891.
\]
Metriky:
\[
\mathrm{precision}=0.90,\qquad \mathrm{recall}=0.81,
\]
\[
\mathrm{accuracy}=\frac{81+9891}{10000}=0.9972,
\]
\[
F_1=\frac{2\cdot0.90\cdot0.81}{0.90+0.81}\approx0.85.
\]

\textbf{Volba v praxi.} Druhý nástroj je lepší v accuracy, precision i $F_1$ a generuje
mnohem méně falešných poplachů. První má vyšší recall, tedy zachytí více anomálií. Pokud
jsou falešně negativní kritické, lze zvolit první nástroj. Pokud je důležitá provozní
zátěž správců a důvěryhodnost hlášení, je vhodnější druhý.
\end{odpoved}

\subsection{2025-podzim, otázka 30: Lineární regrese a L2 regularizace}

Zadání: pro data $(-1,1)$, $(0,3)$, $(1,2)$ určete OLS model $f(x)=\beta_0+\beta_1x$ a
popište vliv L2 regularizace.

\begin{odpoved}
\textbf{Model a OLS.} Lineární regrese s jedním příznakem predikuje
\[
f(x)=\beta_0+\beta_1x.
\]
Metoda nejmenších čtverců minimalizuje
\[
\sum_i(f(x_i)-y_i)^2.
\]
V maticovém zápisu s prvním sloupcem jedniček platí
\[
\beta=(X^TX)^{-1}X^Ty,
\]
pokud je $X^TX$ regulární.

\textbf{Výpočet.} Pro zadaná centrovaná $x$ platí $\sum x_i=0$, takže
\[
\beta_0=\bar y=\frac{1+3+2}{3}=2.
\]
Sklon je
\[
\beta_1=\frac{\sum_i x_i(y_i-\bar y)}{\sum_i x_i^2}
=\frac{(-1)(-1)+0\cdot1+1\cdot0}{1+0+1}
=\frac12.
\]
OLS model je tedy
\[
f(x)=2+\frac12x.
\]
Predikce na $x=-1,0,1$ jsou $1.5$, $2$ a $2.5$.

\textbf{L2 regularizace.} Ridge regrese minimalizuje
\[
\sum_i(f(x_i)-y_i)^2+\alpha\beta_1^2,\qquad \alpha>0,
\]
kde bias $\beta_0$ se typicky nepenalizuje. Regularizace tlačí sklon k nule, takže
$\beta_1$ bude menší než $0.5$, ale zůstane kladný. Proto bude regularizovaná přímka
plošší. V tomto centrovaném příkladu se intercept nezmění a zůstane přibližně $2$;
obecně se intercept může změnit nepřímo podle dat a penalizovaných koeficientů.
\end{odpoved}

\subsection{2025-leto, otázka 4: Centrální limitní věta}

Zadání: formulujte CLV a odhadněte pravděpodobnost, že zpracování 100 vstupů s průměrem
100 ms a směrodatnou odchylkou 20 ms bude trvat víc než 10.3 s.

\begin{odpoved}
\textbf{Znění CLV.} Nechť $X_1,X_2,\ldots$ jsou nezávislé stejně rozdělené náhodné veličiny
se střední hodnotou $\mu$ a rozptylem $\sigma^2>0$. Pak
\[
Y_n=\frac{(X_1+\cdots+X_n)-n\mu}{\sigma\sqrt n}
\]
konverguje v distribuci k normálnímu rozdělení $N(0,1)$. Ekvivalentně, pro distribuční
funkce $F_{Y_n}$ platí $F_{Y_n}(x)\to\Phi(x)$ pro každé reálné $x$.

\textbf{Výpočet.} Měříme v sekundách:
\[
\mu=0.1,\qquad \sigma=0.02,\qquad n=100.
\]
Hledáme
\[
P(X_1+\cdots+X_{100}>10.3).
\]
Standardizace dává
\[
z=\frac{10.3-100\cdot0.1}{0.02\sqrt{100}}
=\frac{0.3}{0.2}=1.5.
\]
Podle CLV
\[
P(X_1+\cdots+X_{100}>10.3)\approx P(Z>1.5)=1-\Phi(1.5).
\]
Z tabulky $\Phi(1.5)\approx0.93$, tedy pravděpodobnost je přibližně
\[
1-0.93=0.07.
\]
Odhad je asi $7\,\%$.

\textbf{Předpoklady.} CLV nepotřebuje normální rozdělení jednotlivých časů, ale potřebuje
nezávislost, stejné rozdělení a konečný rozptyl. Nezávislost může být v praxi problematická:
přetížený počítač zpomalí mnoho běhů současně. Také používáme limitní větu pro $n=100$,
což je aproximace, ale pro běžné aplikace CLV už obvykle rozumná.
\end{odpoved}

\subsection{2024-leto, otázka 30: Metoda podpůrných vektorů}

Zadání: pro lineárně separabilní data tříd $A=\{(1,1),(2,2),(2,0)\}$ a
$B=\{(0,0),(1,0),(0,1)\}$ popište lineární SVM, rozhodovací hranici, vliv bodu $(3,3)$
a použití RBF jádra.

\begin{odpoved}
\textbf{Lineární separabilní SVM.} Pro třídy $t_i\in\{-1,1\}$ hledá SVM nadrovinu
\[
w^Tx+b=0
\]
s maximálním marginem. V kanonickém škálování řeší
\[
\min_{w,b}\frac12\|w\|^2
\quad\text{za podmínek}\quad
t_i(w^Tx_i+b)\ge1.
\]
Šířka marginu je $2/\|w\|$. Rozhodovací hranici určují podpůrné vektory, tedy body ležící
na okrajích marginu.

\textbf{Hranice pro data.} Body třídy $B$ leží vlevo dole a body třídy $A$ vpravo nahoře,
respektive vpravo. Přirozená separace je zhruba podle součtu souřadnic. Hranice může být
přibližně
\[
x_1+x_2=1.5.
\]
Na jedné straně jsou $B$: $(0,0)$ má součet $0$, $(1,0)$ a $(0,1)$ mají součet $1$.
Na druhé straně jsou $A$: $(1,1)$ má součet $2$, $(2,0)$ a $(2,2)$ mají součty $2$ a $4$.
Nejbližší body k hranici jsou zejména $(1,0)$, $(0,1)$, $(1,1)$ a $(2,0)$, takže právě ony
ovlivňují margin.

\textbf{Přidání bodu $(3,3)$.} Bod $(3,3)$ třídy $A$ je daleko od rozhodovací hranice a
neleží na marginu. U hard-margin SVM tedy nebude podpůrným vektorem a hranici prakticky
nezmění. SVM je určeno hlavně nejbližšími body, ne hluboko správně klasifikovanými body.

\textbf{RBF jádro.} Na těchto datech RBF jádro klasifikaci nezlepší podstatně, protože data
jsou už lineárně separabilní a lineární hranice je jednoduchá. RBF jádro pomáhá u nelineárních
struktur, například když jedna třída tvoří malý kruh uprostřed a druhá třída prstenec kolem
ní, nebo u XOR uspořádání. Pak lokální podobnost
\[
K(x,z)=\exp(-\gamma\|x-z\|^2)
\]
umožní nelineární rozhodovací hranici.
\end{odpoved}
